\chapter{四只原主标的：更新后的事件验证逻辑}

\section{恒逸石化：盈利已验证，低位逻辑失效}

恒逸实际H1归母利润CNY59.02亿元，较FY2025的低基数显著改善，说明炼化利润修复并非空泛叙事。但8月10日CNY18.18收盘位于一年区间99.96\%，且高于国信CNY16.50--17.60区间。H1 OCF仍为负CNY4.88亿元。结论是将其置于“盈利和现金验证”，而非将实际利润年化后继续追认低位空间。

\textbf{证据：} F-000703-H1，CNINFO公告1225467439（2026-08-11，第10页）；B-000703，国信证券（2026-07-05，PDF首段/目标区间）。

\section{宏桥控股：低区间价格与利润预告，尚未形成可发布模型}

宏桥处于一年区间30.94\%，Q1 OCF为正CNY92.38亿元，H1预告归母为CNY150--160亿元。招银国际CNY29参照有原始PDF，但其预测采用调整后净利润，发行人H1预告为归母利润，且H1收入和OCF未披露。实际H1及口径调节是第一道升级门槛；在此之前，CNY29仅是外部比较，不是本报告目标。

\textbf{证据：} F-002379-H1，CNINFO公告1225418343（2026-07-11，第1--2页，预告）；B-002379，招银国际（2026-05-13，第1页，13.2倍2026E PE、调整后利润口径）。

\section{亿纬锂能：目标分歧本身就是风险信息}

亿纬处于一年区间22.45\%，H1归母预告为CNY31.30--33.71亿元。东吴CNY103参照建立在30倍2026E PE及不同EPS上，群益CNY77参照未充分披露目标推导；二者不能平均。Q1 OCF为负CNY3.66亿元，实际H1收入与OCF亦未获得。因此本期结论是等待经营与现金数据，而非将两个券商目标的中间值伪装为一致预期。

\textbf{证据：} F-300014-H1，CNINFO公告1225371777（2026-06-15，第1--2页，预告）；B-300014-DW，东吴证券（2026-06-16，30倍2026E PE）；B-300014-QY，群益证券（2026-06-26，第1页，CNY77，目标推导未披露）。

\section{雅化集团：最高外部弹性，不等于最高确定性}

雅化处于一年区间32.67\%，东吴CNY42相对当前CNY19.09有约120\%机械空间。该空间高度依赖单一券商的2026E EPS和锂周期经营假设。公司H1预告归母CNY11--13亿元，而该券商FY2026归母预测为CNY30.32亿元，机械上仍需H2贡献约CNY17.32--19.32亿元；FY2025及Q1 OCF均为负。因此其是最需要跟踪的高弹性验证案例，而不是自动延续的条件核心或当前目标价。

\textbf{证据：} F-002497-H1，CNINFO公告1225411869（2026-07-07，第1--2页，预告）；B-002497，东吴证券（2026-07-08，第1页，16倍2026E PE、CNY42）；FY2025/Q1 OCF见FIN-RAW官方报表包。
